% UNREGISTERED D.72 candidate. Await root and independent proof audits.
\subsection{Compact exact fibres after fixing class two}

\begin{statement}{D.72}{Draft}{Finite nilpotent failure on the requested curve and the Fermat curves}
Let \(U/k\) be a smooth geometrically connected affine hyperbolic
curve over a number field, with a chosen actual arithmetic section
\(s_0\). A rational point or rational tangential cusp supplies such
a choice in every application below. Fix a coefficient prime
\(\ell\), including \(\ell=2\), and use the action induced by \(s_0\).
Put
\[
 P=\pi_1(U_{\bar k})^{(\ell)},\qquad
 T_n=P/\gamma_{n+1}P.
                                                        \tag{D.72.1}
\]
For every fixed actual continuous cocycle
\(\alpha\in Z^1(G_k,T_2)\), with no ramification restriction,
the space
\[
 {\mathscr Z}_n(\alpha)=
   \{z\in Z^1(G_k,T_n):z\bmod\gamma_3=\alpha\},
                 \qquad n\ge2,                         \tag{D.72.2}
\]
is compact in the uniform topology on actual continuous cochains.
More precisely, it is the image of a closed subset of a finite
power of \(\mathbb Z_\ell\) under a continuous rational-polynomial
family of Hall coordinates with coefficients in fixed finite
spaces of continuous cochains.

The corresponding section-class fibres
\[
 {\mathscr S}_n([\alpha])=
   H^1(G_k,T_n)\mathbin{\times}_{H^1(G_k,T_2)}\{[\alpha]\}
                                                        \tag{D.72.3}
\]
are compact Hausdorff. So are the full pro-\(\ell\) fibre and
their exact versions. All transition maps have closed image.
If a fixed finite-stage class lifts separately to every higher
nilpotent stage, it lifts to an actual continuous \(P\)-section.
There is no assumption that the initially supplied lifts agree.
For a once-punctured elliptic curve these compactness and
compatibility assertions also hold for the actual metabelian
group \(P/P''\) and all its nilpotent quotients.

For every generic parameter furnished by D.71 on the requested
once-punctured elliptic curve, a finite-stage class different
from its cusp class fails to lift at some finite nilpotent stage.
The failure is uniform outside any fixed neighborhood of that
cusp in the compact fibre: for \(m\ge2\) and an open neighborhood
\(W\) of its cusp class, some \(N\ge m\) satisfies
\[
  \operatorname{im}\bigl({\mathscr S}_N([\alpha])
            \longrightarrow{\mathscr S}_m([\alpha])\bigr)
                      \subset W.                       \tag{D.72.4}
\]
The same assertion holds on every Fermat curve in D.70, for
its generic parameters and the fixed class-two image of the
chosen cusp. It holds at every coefficient prime. The all-prime
product statement concerns the actual pronilpotent tower; it
has the same compactness and finite-failure conclusions for
the product generic cusp parameters. It does not identify an
arbitrary full profinite section with a pronilpotent section.

No least stage, numerical bound, or first obstruction class is
computed. In particular this does not replace the open explicit
finite-derived obstruction in D.5.
\end{statement}

\begin{proof}
\paragraph{The central lattices and their weights.}
Write \(X\) for the smooth proper compactification, \(D=X-U\),
and \(V=P^{\rm ab}\). The geometric group \(P\) is free
pro-\(\ell\) of rank \(2g+|D(\bar k)|-1\). Its actual
lower-central factors are finite free lattices
\[
 L_j=\gamma_jP/\gamma_{j+1}P
       =\operatorname{Lie}_j(V),\qquad j\ge1,            \tag{D.72.5}
\]
where the right side is the degree-\(j\) part of the free
integral Lie algebra. Hall's basic-commutator theorem and
its pro-prime completion give this identification
\cite[Theorems 5.6 and 6.5, and completion case (iii)]{DE_HallEdmonton},
as checked in D.68. The integral Hall bases apply in every
class and at every prime. Inner automorphisms act trivially
on each \(L_j\); its arithmetic action therefore depends
only on the canonical action on \(V\), regardless of the
ramification of \(s_0\) or a twisting cocycle.

Choose one finite set \(S\) containing the archimedean places,
the places above \(\ell\), and all places of bad reduction
of the marked pair \((X,D)\). At \(v\notin S\), the
canonical action on \(V\) is unramified. Its boundary and
proper parts fit into the integral exact sequence
\[
 0\longrightarrow
 \left(\mathbb Z_\ell[D(\bar k)]/
                   \mathbb Z_\ell{\textstyle\sum_{d\in D(\bar k)}}d\right)(1)
 \longrightarrow V\longrightarrow T_\ell J_X
 \longrightarrow0.                                    \tag{D.72.6}
\]
The exact boundary description and good marked specialization
are the same as in D.70. Arithmetic Frobenius \(F_v\), with
\(q=N(v)\), has boundary eigenvalues \(q\zeta\) for roots
of unity \(\zeta\). Its proper eigenvalues are algebraic
numbers with every complex absolute value \(q^{1/2}\)
\cite[II, Theorem 1.1; IV, Section 3]{fghMilneAV2008}.
This convention agrees with \(F_v(\tau)F_v^{-1}=\tau^q\)
on tame inertia.

After tensoring with \(\mathbb Q_\ell\), \(L_j\) is a
subspace of \(V^{\otimes j}\), or equivalently a
subquotient suffices here. Its Frobenius eigenvalues are
therefore among products of \(j\) eigenvalues of \(V\).
Their complex absolute values belong to
\[
       \{q^{w/2}: j\le w\le2j\}.                       \tag{D.72.7}
\]
No semisimplicity of the extension (D.72.6) is needed:
characteristic polynomials in an invariant filtration
multiply, and tensor-product eigenvalues are the products
of the eigenvalues of its factors. Thus for every \(j\ge3\)
\[
     q\text{ is not an eigenvalue of }F_v\text{ on }
                                   L_j\otimes\mathbb Q_\ell.
                                                        \tag{D.72.8}
\]
The same reasoning excludes eigenvalue \(1\) in every
positive degree, including degrees one and two.

\paragraph{Every central one-cocycle has fixed finite support.}
Let \(j\ge3\) and let \(z:G_k\to L_j\) be any actual
continuous one-cocycle. No finite ramification has been
assumed. At \(v\notin S\), its restriction to inertia
is a homomorphism because the coefficient action there
is trivial. Wild inertia and the prime-to-\(\ell\) tame
factors map trivially to the additive pro-\(\ell\) lattice.
For an \(\ell\)-primary tame generator \(\tau\), put
\(b=z(\tau)\). The actual tame relation and the cocycle
identity give
\[
                       F_vb=q b.                     \tag{D.72.9}
\]
The terms involving \(z(F_v)\) cancel because inertia
acts trivially on the coefficient. Equation (D.72.8)
and torsion-freeness imply \(b=0\). Thus \(z\) kills
every inertia group outside \(S\). The coefficient
action itself is unramified outside \(S\), so the
affine representation \((z,\rho_{L_j})\) kills their
closed normal closure. Consequently
\[
 Z^1(G_k,L_j)=Z^1(G_{k,S},L_j),\qquad j\ge3,             \tag{D.72.10}
\]
as actual continuous cocycle modules.

For each \(r\), the group
\(H^1(G_{k,S},L_j/\ell^rL_j)\) is finite
\cite[I, Corollary 4.15]{DE_MilneADT}. This includes
\(\ell=2\) and real places; only degree-one finiteness
is used, not a cohomological-dimension assertion.
The finite invariant groups satisfy the
Mittag--Leffler condition. Continuous comparison
\cite[Proposition 1.6 and Theorem 2.2]{DE_Jannsen1988}
therefore identifies integral \(H^1\) with the compact
inverse limit of these finite first-cohomology groups.
The multiplication-by-\(\ell\) coefficient sequence
injects its quotient by \(\ell\) into the finite
group \(H^1(G_{k,S},L_j/\ell L_j)\). Compact Nakayama
makes this integral \(H^1\) finitely generated over
\(\mathbb Z_\ell\).

Choose actual continuous representatives of finitely
many module generators, and add the coboundaries of
a lattice basis of \(L_j\). These finitely many
actual maps generate every one-cocycle. The map
from the finite parameter power of \(\mathbb Z_\ell\)
to their linear combinations is continuous in the
uniform topology. Hence
\[
 Z^1(G_k,L_j)\text{ is a finitely generated compact }
              \mathbb Z_\ell\text{-module},\qquad j\ge3.
                                                        \tag{D.72.11}
\]
This argument retains actual coboundaries and works
even if \(H^1\) has finite torsion.

The starting level two is essential for this proof.
In degree two, the proper-proper tensor terms may
have Frobenius eigenvalue \(q\). On a once-punctured
elliptic curve that entire central factor is
\(\bigwedge^2T_\ell E=\mathbb Z_\ell(1)\), whose
unrestricted Kummer group is not finitely generated.
Fixing its actual class-two cocycle removes this
uncontrolled variation. For that curve (D.72.7)
is in fact pure weight \(j\).

\paragraph{Polynomial families and actual coboundaries.}
We give the finite-parameter step to specify fully
why the nonseparated global \(H^2\) creates no gap.
Hall coordinates identify every \(T_n\) homeomorphically
with a finite power of \(\mathbb Z_\ell\). Products,
inverses and powers are rational polynomials with
bounded denominators at fixed \(n\), integral on the
actual coordinate lattice. Appending zero top
coordinates gives a continuous polynomial
set-section \(j_n:T_{n-1}\to T_n\).
The actual action of \(G_k\) on \(T_n\) is polynomial
in the coordinates of its input, with coefficients
given by fixed continuous functions of the Galois
element: substitute the images of the finitely
many free generators in the Hall multiplication
and power polynomials. In particular this is the
actual arithmetic action, with all its higher
coefficients and any ramification of the chosen
reference section retained.

For a finite free lattice \(L\), fix finitely many
continuous rational two-cochains \(\zeta_i\) and put
\[
 B=\left\{u\in\mathbb Z_\ell^a:
          {\textstyle\sum_i}u_i\zeta_i
                       \in dC^1_{\rm cts}(G_k,L)\right\}.
                                                        \tag{D.72.12}
\]
This is the submodule of coefficients giving
\emph{actual integral continuous coboundaries}.
It is a submodule of the Noetherian module
\(\mathbb Z_\ell^a\), hence finite free and closed:
a finite module basis maps a compact finite
power of \(\mathbb Z_\ell\) onto it. Choose one
integral continuous primitive for each basis
vector of \(B\). The basis-coordinate inverse
extends to a \(\mathbb Q_\ell\)-linear map on
\(\mathbb Q_\ell^a\), by a left inverse of its
basis matrix. Applying that fixed linear map
to a polynomial vector taking values in \(B\)
therefore supplies an integral continuous
primitive, polynomial in the parameters and
lying in the finite span of the chosen
primitives. No closure of \(dC^1\) is taken
anywhere in this construction.

At level two, \(\mathscr Z_2(\alpha)=\{\alpha\}\)
is already a finite-parameter family. Suppose
all exact level-\((n-1)\) cocycles are covered
by a family \(f_t\) on a closed set
\(X\subset\mathbb Z_\ell^a\), every coordinate
being a finite sum
\[
             \sum_i P_i(t)\varphi_i(g),
      \quad P_i\in\mathbb Q_\ell[t],\quad
             \varphi_i\in C(G_k,\mathbb Q_\ell).
                                                        \tag{D.72.13}
\]
The set-lifts \(\widehat f_t=j_n(f_t)\) have
central multiplication defects
\[
 \eta_t(g,h)=\widehat f_t(g)\,
     \sigma_0(g)\bigl(\widehat f_t(h)\bigr)\,
                           \widehat f_t(gh)^{-1}
                       \quad\text{in }L_n.             \tag{D.72.14}
\]
Expanding the finitely many Hall polynomials
and the fixed arithmetic coordinates gives
finitely many parameter monomials, each with
a fixed continuous two-cochain as coefficient.
This finite span is uniform over \(X\).
Absorbing the bounded rational denominators
into those cochains, write
\(\eta_t=\sum_i v_i(t)\zeta_i\) with
\(v_i\in\mathbb Z_\ell[t]\).

Exactly the parameters in the closed set
\(X'=X\cap v^{-1}(B)\) lift, with \(B\)
defined by (D.72.12) for the actual lattice
\(L_n\). Its chosen primitive cancels the
defect. Adding arbitrary linear combinations
of the finite generators (D.72.11) gives
every exact lift, since two lifts with the
same projection differ by a central actual
one-cocycle. The new parameter set is
\(X'\times\mathbb Z_\ell^b\), closed and compact.
Its Hall coordinates again have the form
(D.72.13), including after the primitive
choice and every previous lifting step.
This proves the claimed finite-parameter
assertion by induction. Joint continuity
on the compact parameter set and \(G_k\)
gives continuity into the uniform cochain
space, and proves (D.72.2) compact.

\paragraph{Conjugators and compatible inverse limits.}
For any section action on \(T_2\), a fixed
element first projects to \(V^{G_k}\), which
is zero by the absence of eigenvalue one
at a good Frobenius. It then belongs to
\(L_2^{G_k}=0\), by the same argument.
These assertions are integral, since the
lattices are torsion-free. Twisting by
\(\alpha\) does not change the action on
either graded factor. Thus
\[
                       T_2^{G_k,\alpha}=1.             \tag{D.72.15}
\]
Every section class above \([\alpha]\)
has an exact representative above \(\alpha\):
lift its class-two conjugator to \(T_n\).
A conjugator between exact representatives
has fixed class-two image and hence lies
in \(K_n=\ker(T_n\to T_2)\).
Therefore (D.72.3) is the quotient of
the compact Hausdorff space (D.72.2)
by the continuous action of the compact
group \(K_n\). It is compact Hausdorff,
and every transition map has closed image.

Fix an exact representative \(z_m\) of
a class at a finite stage \(m\ge2\).
Its exact lift space at level \(n\ge m\)
is a closed subset of (D.72.2). A class
lift can be made exact by lifting the
required conjugator from \(T_m\), so
individual class liftability makes each
of these compact spaces nonempty.
Every finite compatibility condition
can be met at its highest level and
then projected down. Compactness gives
an inverse system of actual compatible
cocycles. Since every finite quotient
of \(P\) is a nilpotent \(\ell\)-group,
\[
                         P=\varprojlim_n T_n.          \tag{D.72.16}
\]
Their limit is an actual continuous
\(P\)-cocycle. The uniform topology
into \(P\) is exactly the inverse-limit
topology on these exact fibres: each
finite quotient detects a uniform
neighborhood and factors through some
\(T_n\). This proves compactness of
the full exact fibre and, on quotienting
by \(\gamma_3P\), of its class fibre.
Nested nonempty compact conjugator
sets show that two full classes with
identical classes at every level are
already conjugate. Thus no separate
conjugacy ambiguity remains in taking
the inverse limit of section classes.

For a once-punctured elliptic curve,
the metabelian group also has the
needed integral coordinate property.
In rank two the discrete Alexander
module is
\(\mathbb Z[X^{\pm1},Y^{\pm1}]C\),
and the lower-central filtration is
given by powers of \((X-1,Y-1)\),
as used in D.68. Its nilpotent
quotients are torsion-free finitely
generated nilpotent groups. Their
pro-\(\ell\) completions are the
actual metabelian nilpotent quotients,
by the finite-quotient universal
property. The central grade \(j\)
is finite free and, rationally,
a quotient of \(\operatorname{Lie}_j(V)\).
The eigenvalue exclusion (D.72.8)
therefore still applies for \(j\ge3\),
and every preceding polynomial,
cocycle and compactness argument
applies without change. The limit
is the actual metabelian group:
every finite quotient is nilpotent,
and compactness gives surjectivity
onto the inverse limit. This proves
the asserted metabelian compatibility;
it does not assert generic singleton
metabelian fibres on these elliptic
curves.

\paragraph{Finite failure in the named applications.}
For a generic parameter of D.71,
take \(\alpha\) to be its actual
class-two cusp cocycle. The full
pro-\(\ell\) fibre is a singleton.
The intersection of the descending
compact images
\[
 I_{N,m}=
 \operatorname{im}\bigl(\mathscr S_N([\alpha])
                \longrightarrow\mathscr S_m([\alpha])\bigr),
                              \qquad N\ge m,           \tag{D.72.17}
\]
is therefore exactly the cusp class:
a member in every image lifts to
the actual full group by the
compatibility result just proved.
For an open cusp neighborhood \(W\),
if all \(I_{N,m}\setminus W\) were
nonempty, they would be nested
nonempty closed subsets of a
compact space and would have
nonempty intersection. This
contradicts the singleton
intersection. Equation (D.72.4)
and individual finite failure
follow.

The same proof applies to D.70's
Fermat parameters. Their full
pro-\(\ell\) abelian fibre is
already a singleton, so its
subfibre over the specified
class-two cusp is a singleton.
Compactness is now taken after
fixing that class-two cocycle.
No assertion that all elliptic
or Fermat abelian fibres are
compact is used.

Finally the actual full geometric
class-\(n\) quotient is the product
of its pro-prime class-\(n\)
quotients, since finite nilpotent
groups split as products of their
Sylow groups. Cochains, conjugators
and section classes are componentwise.
The product of the compact fibres
just proved is compact, and its
limit is the actual pronilpotent
quotient \(\prod_\ell P_\ell\).
For the product generic cusp
parameters in D.70 or D.71 this
limit fibre is a singleton.
The same compact nested-image
argument gives finite nilpotent
failure and uniform exclusion
outside cusp neighborhoods in
that product tower. This step
concerns the pronilpotent quotient;
the full-profinite cusp conclusions
remain precisely those of D.70
and D.71, using D.24/D.56.
\end{proof}
